% !TeX program = xelatex
% Compile with XeLaTeX in Overleaf.

\documentclass[a4paper]{article}
\usepackage{xeCJK}
\setCJKmainfont[BoldFont=FandolSong-Bold.otf]{FandolSong-Regular.otf}
\usepackage{needspace}

% \usepackage[default]{fontsetup}

\usepackage{fancyhdr}
\usepackage{extramarks}
\usepackage{amsmath}
\usepackage{amsthm}
\usepackage{amsfonts}
\usepackage{tikz}
\usepackage[plain]{algorithm}
\usepackage{algpseudocode}
\usepackage{enumerate}
\usepackage{tikz}

\usetikzlibrary{automata,positioning}

%
% Basic Document Settings
%  

\topmargin=-0.2in
\evensidemargin=0in
\oddsidemargin=0in
\textwidth=6.5in
\textheight=9.5in
\headsep=0.25in

\linespread{1.1}

\pagestyle{fancy}
\lhead{}
\chead{\hmwkClass : \hmwkTitle}
\rhead{\firstxmark}
\lfoot{\lastxmark}
\cfoot{\thepage}

\renewcommand\headrulewidth{0.4pt}
\renewcommand\footrulewidth{0.4pt}

\setlength\parindent{0pt}

%
% Create Problem Sections
%

%修改下标的地方
%\newcommand{\enterProblemHeader}[1]{
	%	\nobreak\extramarks{}{Problem \arabic{#1} continued on next page\ldots}\nobreak{}
	%	\nobreak\extramarks{Problem \arabic{#1} (continued)}{Problem \arabic{#1} continued on next page\ldots}\nobreak{}
	%}

\newcommand{\exitProblemHeader}[1]{
    \nobreak\extramarks{Problem \arabic{#1} }{Problem \arabic{#1} continued on next page\ldots}\nobreak{}
	\stepcounter{#1}
	\nobreak\extramarks{Problem \arabic{#1}}{}\nobreak{}
}

\newcommand*\circled[1]{\tikz[baseline=(char.base)]{
		\node[shape=circle,draw,inner sep=2pt] (char) {#1};}}


\setcounter{secnumdepth}{0}
\newcounter{partCounter}
\newcounter{homeworkProblemCounter}
\setcounter{homeworkProblemCounter}{1}
\nobreak\extramarks{Problem \arabic{homeworkProblemCounter}}{}\nobreak{}

%
% Homework Problem Environment
%
% This environment takes an optional argument. When given, it will adjust the
% problem counter. This is useful for when the problems given for your
% assignment aren't sequential. See the last 3 problems of this template for an
% example.
%

\newenvironment{homeworkProblem}[1][-1]{
	\ifnum#1>0
	\setcounter{homeworkProblemCounter}{#1}
	\fi
	\section{Problem \arabic{homeworkProblemCounter} }
	\setcounter{partCounter}{1}
	% \enterProblemHeader{homeworkProblemCounter} % 注释掉未定义的命令
}{
	\exitProblemHeader{homeworkProblemCounter}
}

%
% Homework Details
%   - Title
%   - Class
%   - Due date
%   - Name
%   - Student ID

\newcommand{\hmwkTitle}{Homework\ \#06}
\newcommand{\hmwkClass}{Probability \& Statistics for EECS}
\newcommand{\hmwkDueDate}{2026-11-10}
\newcommand{\hmwkAuthorName}{杨涵}
\newcommand{\hmwkAuthorID}{2025531064}


%
% Title Page
%

\title{
	\vspace{2in}
	\textmd{\textbf{\hmwkClass:\\  \hmwkTitle}}\\
	\normalsize\vspace{0.1in}\small{Due\ on\ \hmwkDueDate\ at 23:59}\\
	\vspace{4in}
}

\author{
	Name: \textbf{\hmwkAuthorName} \\
	Student ID: \hmwkAuthorID}
\date{}

\renewcommand{\part}[1]{\textbf{\large Part \Alph{partCounter}}\stepcounter{partCounter}\\}

%
% Various Helper Commands
%

% Useful for algorithms
\newcommand{\alg}[1]{\textsc{\bfseries \footnotesize #1}}
% For derivatives
\newcommand{\deriv}[1]{\frac{\mathrm{d}}{\mathrm{d}x} (#1)}
% For partial derivatives
\newcommand{\pderiv}[2]{\frac{\partial}{\partial #1} (#2)}
% Integral dx
\newcommand{\dx}{\mathrm{d}x}
% Alias for the Solution section header
\newcommand{\solution}{\par\addvspace{\baselineskip}\Needspace{4\baselineskip}\noindent\textbf{\large Solution}\par\nobreak\smallskip}
% Probability commands: Expectation, Variance, Covariance, Bias
\newcommand{\E}{\mathrm{E}}
\newcommand{\Var}{\mathrm{Var}}
\newcommand{\Cov}{\mathrm{Cov}}
\newcommand{\Bias}{\mathrm{Bias}}

\begin{document}
	
	
	% \maketitle
	% \thispagestyle{empty}
	% \pagebreak
	
	\date{
		Due on Nov. 10, 2026, 18:00 UTC+8}
	\title{SI 140A-01  Probability \& Statistics for EECS, Fall 2026 \\
		Homework 6}
	\maketitle
	Read all the instructions below carefully before you start working on the assignment, and before you make a submission.
	\begin{itemize}
		\item You are required to write down all the major steps towards making your conclusions; otherwise you may obtain limited points of the problem.
		\item Write your homework in English; otherwise you will get no points of this homework.
		\item Any form of plagiarism will lead to $0$ point of this homework. 
	\end{itemize}
	
	
	
	\newpage
  
	\begin{homeworkProblem}
	\begin{enumerate}[(a)]
\item  For $X \sim \text{Pois}(\lambda)$ and any function $g$, derive the relationship between $E[Xg(X)]$ and $E[g(X+1)]$. This is called the Stein-Chen identity for the Poisson. %(Hint: You can consider the LOTUS.)
\item  Find the moment $E\left(X^{4}\right)$ for $X \sim \operatorname{Pois}(\lambda)$ by using the identity from (a) with the fact that $X$ has mean $\lambda$ and variance $\lambda$.
\end{enumerate}

\solution
\begin{enumerate}[(a)]
    \item The $k=0$ term in the expectation is zero. Starting the sum at $k=1$ and then shifting the index gives
    \begin{align*}
    \E[Xg(X)]
    &= \sum_{k=1}^{\infty} k g(k)e^{-\lambda}\frac{\lambda^k}{k!}\\
    &= \lambda\sum_{k=1}^{\infty}g(k)e^{-\lambda}
       \frac{\lambda^{k-1}}{(k-1)!}\\
    &= \lambda\sum_{k=0}^{\infty}g(k+1)e^{-\lambda}
       \frac{\lambda^k}{k!}\\
    &= \lambda\E[g(X+1)].
    \end{align*}

    \item Since $\E[X]=\lambda$ and $\Var(X)=\lambda$,
    \[
    \E[X^2]=\Var(X)+(\E[X])^2=\lambda+\lambda^2.
    \]
    Applying the identity from (a) with $g(x)=x^2$ gives
    \begin{align*}
    \E[X^3]
    &= \lambda\E[(X+1)^2]\\
    &= \lambda\bigl(\E[X^2]+2\E[X]+1\bigr)\\
    &= \lambda^3+3\lambda^2+\lambda.
    \end{align*}
    Now take $g(x)=x^3$. Then
    \begin{align*}
    \E[X^4]
    &= \lambda\E[(X+1)^3]\\
    &= \lambda\bigl(\E[X^3]+3\E[X^2]+3\E[X]+1\bigr)\\
    &= \lambda^4+6\lambda^3+7\lambda^2+\lambda.
    \end{align*}
\end{enumerate}

\end{homeworkProblem}
	
	\newpage
	\begin{homeworkProblem}
		Alice flips a fair coin $n$ times and Bob flips another fair coin $n+1$ times, resulting in independent $X \sim \operatorname{Bin}(n, \frac{1}{2})$ and $Y \sim \operatorname{Bin}(n+1, \frac{1}{2})$.

(a) Let $V = \min(X,Y)$ be the smaller of $X$ and $Y$, and let $W = \max(X,Y)$ be the larger of $X$ and $Y$. (If $X = Y$, then $V = W = X = Y$.) Find $E(V) + E(W)$.

(b) Show that $P(X < Y) = P(n - X < n + 1 - Y)$.

(c) Compute $P(X < Y)$. Hint: Use (b) and the fact that $X$ and $Y$ are integer-valued.

\solution
\begin{enumerate}[(a)]
    \item Since $V+W=X+Y$, linearity of expectation gives
    \[
    \E[V]+\E[W]=\E[X]+\E[Y]
    =\frac{n}{2}+\frac{n+1}{2}=n+\frac12.
    \]

    \item Swapping heads and tails in every flip produces a sequence with the same probability. Under this swap, the head counts $(X,Y)$ become the tail counts $(n-X,n+1-Y)$. Thus the two pairs have the same joint distribution, and
    \[
    P(X<Y)=P(n-X<n+1-Y).
    \]

    \item Rearranging $n-X<n+1-Y$ gives $Y-X<1$. Since $X$ and $Y$ are integer-valued, this is equivalent to $Y-X\leq 0$, or $Y\leq X$. Hence $P(X<Y)=P(Y\leq X)$ by (b). The two events are complementary, so
    \[
    P(X<Y)=\frac12.
    \]
\end{enumerate}

\end{homeworkProblem}
	
	
 

\newpage
\begin{homeworkProblem} 

Consider two probability distributions
$P=(0.75,0.25)$ and $Q=(0.5,0.5)$
over a binary random variable $X\in\{0,1\}$.
All logarithms are base 2.
\begin{enumerate}
    \item[(a)] \textbf{Basic Calculations (5 points)}

    Compute the total variation distance and
    KL divergence:
    \begin{align*}
        d_{\mathrm{TV}}(P,Q)
        &= \frac{1}{2}\sum_x |P(x)-Q(x)|,\\
        D_{\mathrm{KL}}(P\|Q)
        &= \sum_x P(x)\log_2\frac{P(x)}{Q(x)}.
    \end{align*}
    Is KL divergence symmetric?
    Justify your answer.

    \item[(b)] \textbf{Connection to Information Theory
    (5 points)}

    \begin{enumerate}[(i)]
    \item Compute the entropy $H(P)$ and cross-entropy $H(P,Q)$, which are defined as
    \begin{align*}
        H(P) &= -\sum_x P(x)\log_2 P(x),\\
        H(P,Q) &= -\sum_x P(x)\log_2 Q(x).
    \end{align*}
    \item Show that
    \[
        H(P,Q)-H(P)=D_{\mathrm{KL}}(P\|Q).
    \]
    Briefly explain the significance of this
    identity in lossless data compression.
    \end{enumerate}

    \item[(c)] \textbf{Connection to Machine Learning
    (5 points)}

    Suppose $P$ represents the true label
    distribution in a binary classification task.
    Two models predict $Q=(0.5,0.5)$ and
    $R=(0.95,0.05)$, respectively.

    Compute $D_{\mathrm{KL}}(P\|Q)$ and
    $D_{\mathrm{KL}}(P\|R)$.
    Which model has a smaller expected
    cross-entropy loss? Explain your answer.
\end{enumerate}

\solution
\begin{enumerate}[(a)]
    \item The total variation distance is
    \[
    d_{\mathrm{TV}}(P,Q)=\frac12\sum_x |P(x)-Q(x)|
    =\frac12\bigl(|0.75-0.5|+|0.25-0.5|\bigr)=0.25.
    \]
    The KL divergences in the two directions are
    \begin{align*}
    D_{\mathrm{KL}}(P\|Q)
    &= 0.75\log_2\frac{0.75}{0.5}
       +0.25\log_2\frac{0.25}{0.5}
       \approx 0.1887\ \text{bits},\\
    D_{\mathrm{KL}}(Q\|P)
    &= 0.5\log_2\frac{0.5}{0.75}
       +0.5\log_2\frac{0.5}{0.25}
       \approx 0.2075\ \text{bits}.
    \end{align*}
    These values differ, so KL divergence is not symmetric.

    \Needspace{10\baselineskip}
    \item
    \begin{enumerate}[(i)]
        \item The entropy and cross-entropy are
        \begin{align*}
        H(P)
        &= -0.75\log_2 0.75-0.25\log_2 0.25
           \approx 0.8113\ \text{bits},\\
        H(P,Q)
        &= -0.75\log_2 0.5-0.25\log_2 0.5
           =1\ \text{bit}.
        \end{align*}

        \Needspace{9\baselineskip}
        \item Subtracting the entropy from the cross-entropy gives
        \begin{align*}
        H(P,Q)-H(P)
        &= -\sum_x P(x)\log_2 Q(x)+\sum_x P(x)\log_2 P(x)\\
        &= \sum_x P(x)\log_2\frac{P(x)}{Q(x)}
        =D_{\mathrm{KL}}(P\|Q).
        \end{align*}
        In lossless data compression, $D_{\mathrm{KL}}(P\|Q)$ is the extra number of bits per symbol, on average, for ideal coding based on $Q$ when the true source follows $P$, compared with ideal coding based on $P$.
    \end{enumerate}

    \Needspace{12\baselineskip}
    \item From (a), $D_{\mathrm{KL}}(P\|Q)\approx 0.1887$ bits. For $R$,
    \[
    D_{\mathrm{KL}}(P\|R)
    =0.75\log_2\frac{0.75}{0.95}
     +0.25\log_2\frac{0.25}{0.05}
    \approx 0.3247\ \text{bits}.
    \]
    For a model distribution $M$, the expected cross-entropy loss is
    $H(P,M)=H(P)+D_{\mathrm{KL}}(P\|M)$. Since $H(P,Q)=1$ bit and
    $H(P,R)\approx 1.136$ bits, the model predicting $Q$ has the smaller expected cross-entropy loss.
\end{enumerate}

\end{homeworkProblem}

	
\end{document}
